\subsection{NTK recursion}
\begin{theorem}[NTK recursion with trainable biases and scaling]\label{thm:ntk_recursion_scaled}
Assumptions:
\begin{itemize}
    \item NTK parameterization: each layer \( \ell \) has width \( n_\ell \) and scaling \( 1/\sqrt{n_\ell} \).
    \item Sequential infinite-width limit: \( n_1\to\infty,\dots,n_L\to\infty \) with \( w^{(\ell)}_j \sim \mathcal{N}(0,\mathbf{I}) \) and \( b^{(\ell)}_j \sim \mathcal{N}(0,1) \) i.i.d., independent across layers.
    \item Trainable parameters: \( \{A^{(\ell)}_j\}_{j=1}^{n_\ell} \) and \( c^{(\ell)} \) are trained; \( \{w^{(\ell)}_j,b^{(\ell)}_j\} \) are frozen.
    \item Activation: \( \sigma \) is piecewise \( C^1 \) (e.g., ReLU); write \( \dot{\sigma} \) for its derivative a.e. Bias scale \( \beta \in \mathbb{R} \), output scales \( \sigma_A,\sigma_c\in\mathbb{R}_+ \).
\end{itemize}
Then the Neural Tangent Kernel (NTK) satisfies, for all \( \ell\ge1 \) and \( x,x' \in \mathbb{R}^{d_0} \),
\[
\Theta^{(0)}(x,x') \;=\; 0,\qquad
\boxed{\;\Theta^{(\ell)}(x,x') \;=\; \Theta^{(\ell-1)}(x,x')\cdot \Big(\sigma_A^2\,\dot{\Sigma}^{(\ell)}(x,x')\Big) \;+\; \Sigma^{(\ell)}(x,x') \;+\; 1\;}
\]
where (expectations taken over layer-\( \ell \) parameters \( (w,b) \))
\[
\Sigma^{(\ell)}(x,x') \;=\; \mathbb{E}_{w,b}\!\left[\sigma\!\big(w^\top h^{(\ell-1)}(x) + \beta b\big)\,\sigma\!\big(w^\top h^{(\ell-1)}(x') + \beta b\big)\right],
\]
\[
\dot{\Sigma}^{(\ell)}(x,x') \;=\; \mathbb{E}_{w,b}\!\left[\dot{\sigma}\!\big(w^\top h^{(\ell-1)}(x) + \beta b\big)\,\dot{\sigma}\!\big(w^\top h^{(\ell-1)}(x') + \beta b\big)\;\underbrace{w^\top w}_{\text{explicit weight-norm factor}}\right].
\]
Here the additive \( +1 \) is the contribution of the trainable bias \( c^{(\ell)} \) when \( \sigma_c^2=1 \). Equivalently, in fully parameterized form,
\[
\Theta^{(\ell)} \;=\; \Theta^{(\ell-1)}\cdot (\sigma_A^2\,\dot{\Sigma}^{(\ell)}) \;+\; \sigma_A^2\,\Sigma^{(\ell)} \;+\; \sigma_c^2,
\]
and absorbing \( \sigma_A^2 \) into \( \Sigma^{(\ell)} \) yields the boxed formula above. Proof: See Appendix, “Proof of the NTK Recurrence Relation.”
\end{theorem}

\medskip
\noindent\textbf{Corollary (Two-layer NTK).}\label{cor:two_layer_ntk}
Under the assumptions of Theorem~\ref{thm:ntk_recursion_scaled} with \( L=2 \),
\[
\Theta^{(1)}(x,x') \;=\; \Sigma^{(1)}(x,x') \;+\; 1,
\]
\[
\Theta^{(2)}(x,x') \;=\; \big(\Sigma^{(1)}(x,x')+1\big)\cdot \big(\sigma_A^2\,\dot{\Sigma}^{(2)}(x,x')\big) \;+\; \Sigma^{(2)}(x,x') \;+\; 1,
\]
where \( \Sigma^{(1)}(x,x')=\mathbb{E}_{w,b}[\sigma(w^\top x+\beta b)\sigma(w^\top x'+\beta b)] \) and
\[
\dot{\Sigma}^{(2)}(x,x')=\mathbb{E}_{w,b}\!\left[\dot{\sigma}\!\big(w^\top h^{(1)}(x)+\beta b\big)\,\dot{\sigma}\!\big(w^\top h^{(1)}(x')+\beta b\big)\;w^\top w\right].
\]
Proof: Immediate from Theorem~\ref{thm:ntk_recursion_scaled}.

\medskip
\noindent\textbf{Corollary (General L-layer NTK; closed form).}\label{cor:general_L_layer_ntk}
Under the assumptions of Theorem~\ref{thm:ntk_recursion_scaled}, for any \( L\ge1 \),
\[
\boxed{\;\Theta^{(L)}(x,x') \;=\; \sum_{\ell=1}^{L} \big(\Sigma^{(\ell)}(x,x')+1\big)\,\prod_{k=\ell+1}^{L}\Big(\sigma_A^2\,\dot{\Sigma}^{(k)}(x,x')\Big)\;}
\]
with the convention that an empty product equals \( 1 \). If \( \sigma_c^2=0 \) (no trainable bias), the \( +1 \) terms vanish and the expression reduces to the standard sum \( \sum_{\ell=1}^{L} \Sigma^{(\ell)}\prod_{k=\ell+1}^{L}\dot{\Sigma}^{(k)} \). Proof: By expanding the recursion in Theorem~\ref{thm:ntk_recursion_scaled}.

\paragraph{Homogeneous activation hypothesis and two-layer compact form.}
Assume \( \sigma \) is positively 1-homogeneous (e.g., ReLU), so \( \sigma(\alpha z)=\alpha\,\sigma(z) \) for \( \alpha\ge0 \). Then base kernels factor into a radial part and a zonal (angular) part. In the two-layer low-rank setting, if we denote
\[
\rho_1 \;=\; \cos\angle(x_1,y_1),\qquad w_r \;=\; \frac{\|x_1\|\,\|y_1\|}{r},
\]
the NTK admits the compact representation
\[
\Theta^{(2)}(x,y) \;=\; \Psi(\rho_1,w_r),\qquad
\Psi(u,w) \;=\; \frac{2}{r}\,\Theta^{(1)}(u)\!\left(1-\frac{\arccos(u)}{\pi}\right) \;+\; \frac{2}{r}\,w\,\Sigma^{(1)}(u) \;+\; 1,
\]
with \( \Sigma^{(1)}(u)=\frac{1}{\pi}\big(\sqrt{1-u^2}+u(1-\arccos u)\big) \). By rotational invariance (Fisher–Kibble decoupling),
\[
\rho_1 \sim \text{Fisher}(\rho,r),\qquad (\|x_1\|^2,\|y_1\|^2)\sim \text{Kibble}(r,\rho),
\]
and \( \rho_1 \) is independent of \( (\|x_1\|,\|y_1\|) \) to leading order, justifying the radial–angular separation captured by \( \Psi \).
